%% THERMODYNAMIQUE %%
%
\section{Humidit\'e relative \`a partir de la sonde carbone}

\itemstyle{ Op\'erateur~: } {\sf HumCarb}
\index{HumCarb}

%%%%%OBJET%%%%%
\itemstyle{ Objet~: }
Cet op\'erateur calcule l'humidit\'e \`a partir de la tension
de la sonde carbone. Il utilise 3 relations diff\'erentes en 
fonction de la temp\'erature (seuils \`a $-10$ et$+12^\circ$C).\newline
Une correction de vitesse est ensuite appliqu\'ee.

%%%%%ARGUMENTS%%%%%
\bigskip
%%%%%ENTREES%%%%/
\begin{tabarg}
\multicolumn{3}{l}{\titretab{Entr\'ees}}\\
\hline
$V$ & Vecteur, \texttt{float} & Tension de la sonde carbone ($V$) \\
$T_s$ & Vecteur, \texttt{float} & Temp\'erature statique de l'air ($K$ ou $^\circ C$) \\
$T_i$ & Vecteur, \texttt{float} & Temp\'erature d'impact ($K$ ou $^\circ C$) \\
$P_s$ & Vecteur, \texttt{float} & Pression statique ($hPa$)\\
$\Delta P$ & Vecteur, \texttt{float} & Pression dynamique ($hPa$)\\
$C_1$ & Coeff., \texttt{float[4]} & Coefficients pour la gamme : $ T_s< -10^\circ C $\\
$C_2$ &  Coeff., \texttt{float[4]} & Coefficients pour la gamme : $ -10^\circ C \leq T_s \leq +12^\circ C $\\
$C_3$ &  Coeff., \texttt{float[4]} & Coefficients pour la gamme : $ +12^\circ C < T_s $\\
\end{tabarg}

\medskip
%%%%%SORTIES%%%%
\begin{tabarg}
\multicolumn{3}{l}{\titretab{Sorties}}\\
\hline
$H_u$ & Vecteur, \texttt{float} & Humidit\'e relative ($\%$)\\
\end{tabarg}


%%%%%FORMULE%%%%
\itemstyle{ Formule }

% calcul de l'humidite carbone
\[
Hu = \frac{ P_s . ew(T_i)}{(P_s + \Delta P) ew(T_s)}
     .\left( C_m[0] + C_m[1].V + C_m[2].V^2 + C_m[3].V^3 + C_m[4].V^4 \right)
\]
$C_m$ est \'egal \`a $C_1$, $C_2$, ou $C_3$ selon la gamme de 
temp\'erature.\newline
$ew()$ est la tension de vapeur saturante par rapport \`a l'eau.
Si $T_s < 0$, $ew()$ est remplac\'ee par la tension de vapeur saturante
par rapport \`a la glace $ewi()$. 


%%%%COMMENTAIRE%%%%
\itemstyle{Commentaire~:}

Les fonctions $ew()$ et $ewi()$ sont d\'ecrites en annexe \ref{annexe-thermo}. 

%%%%%AUTEUR%%%%%
\itemstyle{ Auteur~:}
CNRM/GMEI/TRAMM

%%%%%REFERENCES%%%%%
\itemstyle{ R\'ef\'erences~:}
Note du CAM sur les appareils de mesure de l'hygrom\'etrie \cite{CAM-hum}
